\section{Métodos}
\subsection{Diseño, datos y alcance del inventario}
Se realizó una reconstrucción matemática de estados de equilibrio. Se adoptó 50 MeV como referencia bibliográfica aproximada y se evaluaron puntos entre 37 y 50 MeV para localizar el cruce de abundancias. Los puntos intermedios fueron evaluaciones numéricas, sin asignarles nuevas épocas físicas. Se utilizaron unidades naturales, $\hbar=c=k_B=1$, con masas y temperaturas en GeV.

Se leyeron las listas PDG2014, PDG2020 y PDG2025 del repositorio Thermal-FIST, fijadas en la revisión \texttt{4086191e9522} \citep{TFLists2026}. El suplemento contiene el identificador completo, archivos originales, licencia y huellas SHA-256. Se conservaron las masas nominales y degeneraciones de espín; se excluyeron líneas comentadas y entradas de degeneración cero. Cada estado de carga se contó según su entrada, y las antipartículas se incluyeron por separado cuando correspondía. Todas las entradas bariónicas activas tenían $|B_i|=1$.

\begin{table}[t]
\centering\small
\caption{Dimensiones de los inventarios utilizados. Las especies activas incluyen antipartículas; $B+\bar B$ cuenta las entradas bariónicas activas.}
\label{tab:inventory}
\setlength{\tabcolsep}{4pt}
\begin{tabular}{lrrr}\toprule
Lista & Filas & Especies activas & $B+\bar B$\\\midrule
PDG2014 & 219 & 387 & 260\\
PDG2020 & 242 & 433 & 298\\
PDG2025 & 248 & 443 & 292\\
\bottomrule\end{tabular}
\end{table}

Las listas abarcan el sector de quarks $u,d,s$; no se agregaron núcleos ni listas de otros sectores de quarks. Se adoptó un HRG ideal, con masas fijas, anchos de resonancia nulos y temperatura común. No se incorporaron correcciones de volumen excluido ni una evolución de reacciones. Se mantuvieron $\mu_Q=\mu_S=0$ y se resolvió $\mu_B$. Esta simplificación no impone neutralidad eléctrica ni extrañeza neta nula; tampoco se resolvieron potenciales químicos leptónicos que compensaran la carga.

\subsection{Densidades y entropía por especie de partícula}
Para cada especie de partícula masiva se calculó
\begin{equation}
n_i=\frac{g_i m_i^2T}{2\pi^2}
\sum_{k=1}^{K_{\max}}\frac{c_{i,k}}{k}
K_2(kx_i)e^{k\alpha_i},\label{eq:number}
\end{equation}
con $x_i=m_i/T$, $\alpha_i=\mu_i/T$ y $\mu_i=B_i\mu_B$. Aquí $g_i=2J_i+1$; $c_{i,k}=1$ para bosones y $c_{i,k}=(-1)^{k+1}$ para fermiones. $K_\nu$ denota la función de Bessel modificada de segunda clase. Estas expresiones implementan las distribuciones de Bose-Einstein y Fermi-Dirac; no describen una prueba estadística sobre mediciones.

La entropía se obtuvo con el mismo potencial químico:
\begin{align}
s_i&=\frac{g_i m_i^2T}{2\pi^2}
\sum_{k=1}^{K_{\max}}c_{i,k}e^{k\alpha_i}\bigg[
\frac{4K_2(kx_i)}{k^2}\nonumber\\
&\qquad+\frac{x_iK_1(kx_i)}{k}
-\frac{\alpha_iK_2(kx_i)}{k}\bigg].\label{eq:entropy}
\end{align}
Se usaron 96 términos en las sumas y 192 en el control de convergencia. Los mesones tuvieron $\mu_i=0$. La entropía total incluyó hadrones, fotones, electrones y positrones, muones y antimuones, y tres sabores de neutrinos y antineutrinos:
\begin{equation}
s_{\rm tot}=s_{\rm had}+s_\gamma+s_{e^\pm}
+s_{\mu^\pm}+s_{\nu\bar\nu}.\label{eq:stotal}
\end{equation}
Fotones, electrones y neutrinos se trataron como componentes sin masa, con degeneraciones conjuntas 2, 4 y 6. Se utilizó $s_{\rm rel}=(2\pi^2/45)gT^3$, multiplicado por $7/8$ para fermiones. Los muones se calcularon como partículas masivas con $m_\mu=0.1056583755$ GeV y degeneración conjunta 4 \citep{PDG2025}. Se fijaron a cero los potenciales químicos del sector ligero.

Se definieron $n_\Sigma=n_B+n_{\bar B}$ y $\Delta n_B=n_B-n_{\bar B}$, junto con $Y_i=n_i/s_{\rm tot}$. En unidades físicas, el denominador del rendimiento es $s_{\rm tot}/k_B$, expresado en las mismas unidades de volumen que $n_i$. Las densidades naturales se convirtieron a cm$^{-3}$ con $\hbar c=1.973269804\times10^{-14}$ GeV cm \citep{Robinson2025}. La conversión de temperatura utilizó $k_B=8.617333262\times10^{-5}$ eV K$^{-1}$.

\subsection{Restricción cosmológica y solución inversa}
Se adoptó $\eta=(6.12\pm0.04)\times10^{-10}$ como entrada externa CMB, compilada por \citet{Fields2025}. La conservación adiabática de entropía y del número bariónico neto dio
\begin{equation}
\frac{\Delta n_B(T,\alpha)}{s_{\rm tot}(T,\alpha)}
=Y_{\Delta B}^{\rm ref}=\frac{\eta}{7.04},
\quad\alpha=\frac{\mu_B}{T}.\label{eq:inverse}
\end{equation}
El factor 7.04 convierte la abundancia por fotón del presente en abundancia por entropía, bajo la aproximación estándar posterior a la aniquilación $e^+e^-$. No se impuso $s(T)/n_\gamma(T)=7.04$ al plasma de 50 MeV.

Para evitar la pérdida de precisión por restar poblaciones casi iguales, se definió $a_{i,k}=g_i m_i^2T c_{i,k}K_2(kx_i)/(2\pi^2k)$ para las entradas $B_i=+1$ y se evaluó
\begin{align}
n_\Sigma&=2\sum_{i,k}a_{i,k}\cosh(k\alpha),\nonumber\\
\Delta n_B&=2\sum_{i,k}a_{i,k}\sinh(k\alpha).\label{eq:stable}
\end{align}
La entropía del par se sumó de forma equivalente mediante funciones hiperbólicas. Se resolvió la ecuación~\eqref{eq:inverse} con el método de Brent, en $0\leq\alpha\leq1$, con tolerancias absoluta $10^{-18}$ y relativa $10^{-13}$. Se recuperaron $n_B=(n_\Sigma+\Delta n_B)/2$ y $n_{\bar B}=(n_\Sigma-\Delta n_B)/2$.

Para contrastar el mismo criterio de \citet{Yang2022} se definió
\begin{equation}
R(T)=\frac{n_{\bar B}}{\Delta n_B}
=\frac{Y_{\bar B}}{Y_{\Delta B}}.\label{eq:ratio}
\end{equation}
Un segundo cálculo de raíz localizó $R(T_\times)=1$ dentro de 37--40 MeV, con tolerancias absoluta $10^{-10}$ MeV y relativa $10^{-13}$. Este intervalo fue un encuadre numérico del cruce, no una temperatura inicial ajustada para producir concordancia. Se repitió la solución con los tres catálogos y con los extremos $\eta=6.08$ y $6.16\times10^{-10}$, variando sólo esa entrada.

\subsection{Comparación con el presente y prueba de unicidad}
Con $T_{\gamma0}=2.7255\pm0.0006$ K, redondeado de \citet{Fixsen2009}, se calcularon
\begin{align}
n_{\gamma0}&=\frac{2\zeta(3)}{\pi^2}
\left(\frac{k_BT_{\gamma0}}{\hbar c}\right)^3,\nonumber\\
s_0/k_B&\simeq7.04n_{\gamma0},\qquad
n_{b0}=\eta n_{\gamma0}.\label{eq:present}
\end{align}
Para comparar cantidades en la misma región comóvil se utilizó $s(T)V(T)=s_0V_0$ y, por tanto,
\begin{equation}
N_i(T)=Y_i(T)(s_0/k_B)V_0.\label{eq:comoving}
\end{equation}
Se eligió $V_0=1$ m$^3$ como normalización explícita; los números resultantes son cantidades medias y pueden ser fraccionarios. No se adoptó un volumen para todo el Universo.

La composición primordial equivalente a la densidad media actual se convirtió usando $Y_p=0.245\pm0.003$ y $(D/H)_{\rm P}=(2.508\pm0.029)\times10^{-5}$ \citep{Fields2025}: $n_{H0}^{\rm prim}\simeq(1-Y_p)n_{b0}$, $n_{{}^4He,0}^{\rm prim}\simeq Y_p n_{b0}/4$ y $n_{D0}^{\rm prim}\simeq(D/H)_{\rm P}n_{H0}^{\rm prim}$. Se despreciaron las correcciones de núcleos traza al balance H/He, las correcciones de enlace, el procesamiento estelar y la ionización. Se propagaron incertidumbres a primer orden suponiendo entradas independientes; las salidas comparten entradas y están correlacionadas.

La unicidad se examinó con el mapa $(Y_B,Y_{\bar B})\mapsto Y_B-Y_{\bar B}$. Su matriz $[1,-1]$ tiene rango uno y dirección nula $(1,1)$. Esta propiedad identifica la información que falta en el residuo, sin proponer historias cinéticas adicionales.

\subsection{Verificación y tratamiento de incertidumbres}
Los 24 controles registrados compararon 96 y 192 términos, comprobaron el cierre de la restricción cosmológica y contrastaron número, exceso y entropía de protones y neutrones con integrales independientes a potencial químico finito. Se verificaron las huellas de los catálogos y se exportaron resultados agregados y por especie de partícula. Las tablas y figuras se generaron directamente de esos archivos.

No se hicieron pruebas de hipótesis, ajustes a observaciones nuevas ni simulaciones aleatorias. La variación entre catálogos es una sensibilidad de inventario y el intervalo por $\eta$ propaga una sola entrada; ninguno representa el error físico total del HRG. Los errores de convergencia miden la implementación numérica. No se calcularon valores $p$ ni una significación estadística de la proximidad a 38.2 MeV.
